\section{Помехоустойчивое кодирование. Кодирование с исправлением ошибок. Возможность исправления всех ошибок.}

\subsection*{Основные понятия}

\paragraph{Помехоустойчивое кодирование} — это способ представления данных, позволяющий обнаруживать или исправлять ошибки, возникшие при передаче по каналу с шумом.
Главная идея заключается во введении **избыточности**: из всех возможных кодовых слов $2^n$ мы используем только некоторое подмножество «разрешенных» слов. Если принятое слово не входит в это подмножество, значит, произошла ошибка.

\subsection*{Расстояние Хэмминга}

Для анализа ошибок Новиков использует метрику в пространстве двоичных слов.
\paragraph{Расстояние Хэмминга $d(x, y)$} — это количество позиций (разрядов), в которых слова $x$ и $y$ одинаковой длины различаются.
\paragraph{Минимальное расстояние кода $d_{min}$} — это минимальное из расстояний Хэмминга между всеми парами различных разрешенных кодовых слов.

\subsection*{Обнаружение и исправление ошибок}

Способность кода бороться с помехами напрямую зависит от $d_{min}$:

\begin{enumerate}
    \item \textbf{Обнаружение $k$ ошибок:} 
    Чтобы гарантированно заметить $k$ ошибок, необходимо, чтобы $d_{min} \ge k + 1$. 
    \textit{Пример:} Код с проверкой на четность ($d_{min}=2$) обнаруживает 1 ошибку, но не может её исправить.
    
    \item \textbf{Исправление $k$ ошибок:} 
    Чтобы гарантированно исправить $k$ ошибок, необходимо, чтобы $d_{min} \ge 2k + 1$.
    \textit{Логика:} Испорченное слово должно остаться «ближе» к исходному разрешенному слову, чем к любому другому. Это называется \textbf{декодированием по принципу максимального правдоподобия}.
\end{enumerate}

\subsection*{Геометрическая интерпретация: Сферы Хэмминга}
Вокруг каждого разрешенного кодового слова можно представить «сферу» радиуса $k$. Если сферы радиуса $k$ вокруг разных кодовых слов не пересекаются, то любую ошибку кратности $\le k$ можно исправить, «притянув» испорченное слово обратно к центру сферы.

\subsection*{Возможность исправления всех ошибок}

\paragraph{Ограничения}
Можно ли исправить \textbf{вообще все} ошибки? 
\begin{itemize}
    \item \textbf{Нет}, если количество ошибок превышает исправляющую способность кода ($k > \lfloor \frac{d_{min}-1}{2} \rfloor$). В этом случае слово может «улететь» в сферу другого кодового слова, и система совершит ложное исправление.
    \item \textbf{Теорема Шеннона:} Для канала с шумом существует такая скорость передачи данных $R < C$ (пропускная способность), при которой вероятность ошибки можно сделать сколь угодно малой, но не равной нулю.
\end{itemize}

\paragraph{Совершенные коды}
Это коды, в которых сферы Хэмминга не только не пересекаются, но и \textbf{плотно упаковывают} все пространство (каждое возможное слово либо является разрешенным, либо находится в сфере радиуса $k$ ровно одного разрешенного слова). Пример: Коды Хэмминга.

\subsection*{Примеры}
\begin{itemize}
    \item \textbf{Код с повторением (3-кратным):} $0 \to 000, 1 \to 111$. 
    $d_{min}=3$. Можно обнаружить 2 ошибки или исправить 1 ошибку (мажоритарный метод: «голосование» большинством).
    \item \textbf{Код Хэмминга (7, 4):} К 4 информационным битам добавляется 3 контрольных. Позволяет исправлять любую одиночную ошибку.
\end{itemize}